\documentclass[a4paper,openright,10pt]{article}
\usepackage{amsmath}
\usepackage{amssymb}
\usepackage{graphicx}
\usepackage{fancyhdr}
\pagestyle{fancy}
\fancyhf{}
\rfoot{P\'agina \thepage \hspace{1pt} de 2}
\everymath{\displaystyle}

\begin{document}

\title{Probabilidad y Estadística}
\author{Prof. Mtro. Luis Jos\'{e} Yudico Anaya}
\date{14-10-2024}
\maketitle

\begin{center}
\textbf{INGENIERÍA EN NANOTECNOLOGÍA}\\
\vspace{.5cm}
\textbf{2024-2025-I}\\
\vspace{.5cm}
\textbf{EJERCICIOS SOBRE PROBABILIDAD BÁSICA}
\end{center}


\section*{Ejercicio 1: Espacios Muestrales}
Un experimento consiste en lanzar un dado de seis caras. Define el espacio muestral y enumera todos los eventos posibles que se pueden obtener al lanzar el dado.

\section*{Ejercicio 2: Eventos}
En una baraja de cartas estándar (52 cartas), define el evento \( A \) de sacar una carta roja y el evento \( B \) de sacar un número par. ¿Cuál es la probabilidad de que ocurra el evento \( A \cup B \)?

\section*{Ejercicio 3: Experimentos con Resultados Igualmente Probables}
Si se lanzan dos monedas, ¿cuál es la probabilidad de que al menos una de ellas muestre cara? Usa el espacio muestral para calcularlo.

\section*{Ejercicio 4: Funciones de Masa de Probabilidad}
Un dado cargado tiene las siguientes probabilidades de resultados: \( P(1) = 0.1 \), \( P(2) = 0.2 \), \( P(3) = 0.3 \), \( P(4) = 0.2 \), \( P(5) = 0.1 \), \( P(6) = 0.1 \). ¿Cuál es la probabilidad de obtener un número mayor que 4?

\section*{Ejercicio 5: Axiomas de Probabilidad}
Demuestra que para cualquier evento \( A \) en un espacio muestral \( S \), se cumple que \( P(A) + P(A') = 1 \).

\section*{Ejercicio 6.1: Probabilidad Condicional}
En un examen, el 80\% de los estudiantes estudian para el examen. De los que estudian, el 90\% aprueba el examen, mientras que el 50\% de los que no estudian aprueban. ¿Cuál es la probabilidad de que un estudiante haya estudiado dado que aprobó el examen?


\section*{Ejercicio 6.2: Probabilidad Condicional}
Se lanzan dos dados. ¿Cuál es la probabilidad de que la suma de los dos dados sea 8 dado que el primer dado muestra un 5?

\section*{Ejercicio 7: Teorema de Bayes}
En una fábrica, el 90\% de los productos son de calidad alta. Un 5\% de los productos de calidad alta son defectuosos y el 20\% de los productos de calidad baja son defectuosos. Si se selecciona un producto defectuoso, ¿cuál es la probabilidad de que sea de calidad alta?

\section*{Ejercicio 8: Eventos Independientes}
Si se lanza una moneda y se tira un dado, ¿son independientes los eventos "salir cara en la moneda" y "sacar un número impar en el dado"? Justifica tu respuesta.

\section*{Ejercicio 9: Regla de Multiplicación}
En un examen, la probabilidad de que un estudiante apruebe la primera parte es de 0.7 y de que apruebe la segunda parte es de 0.8. ¿Cuál es la probabilidad de que el estudiante apruebe ambas partes, suponiendo que los eventos son independientes?

\section*{Ejercicio 10: Aplicaciones de la Probabilidad}
Un sistema tiene tres componentes, y cada componente tiene una probabilidad de fallo de 0.1. Si el sistema falla si al menos uno de los componentes falla, ¿cuál es la probabilidad de que el sistema funcione correctamente?

\end{document}
